\documentclass[11pt]{article}
\usepackage[margin=1in]{geometry}
\usepackage{amsmath,amssymb,amsthm}
\usepackage[round]{natbib}
\usepackage{hyperref}
\usepackage{microtype}

\newtheorem{theorem}{Theorem}
\newtheorem{proposition}{Proposition}
\newtheorem{lemma}{Lemma}
\newtheorem{corollary}{Corollary}
\newtheorem{conjecture}{Conjecture}
\newtheorem{remark}{Remark}

\title{The Width of the Conformal Fan:\\
Dependence and the Variance of Realized Coverage}
\author{Peter Cotton}
\date{Draft --- companion note}

\begin{document}
\maketitle

\begin{abstract}
Split conformal prediction's realized, calibration-conditional coverage fluctuates around the
nominal $1-\alpha$. For exchangeable-but-independent calibration scores it follows the
$\mathrm{Beta}(k,n-k+1)$ law, variance about $\alpha(1-\alpha)/n$, the classical fan. We show the
width of that fan is governed by the sign of the cross-sample dependence. The mean stays pinned at
the marginal level whatever the dependence; the variance, to leading order, is the average pairwise
covariance of the exceedance indicators at the operating quantile times the independent fan.
Positive (extendable) dependence adds a between-dataset term and widens the fan; negative dependence
narrows it, to exactly zero at the maximally negatively associated contest floor
$\rho=-1/(n-1)$, where the realized coverage equals the nominal level identically. We prove the
positive side exactly through de Finetti and the law of total variance, the negative side exactly at
the floor and to leading order in general, and reduce the remaining finite-sample inequality to a
dispersive ordering supported numerically across dependence structures. The governing sign is the
one the companion note attaches to the finite de Finetti measure: a genuine prior widens the fan, the
signed corner collapses it.
\end{abstract}

\section{Setup}

Split conformal returns a set with marginal coverage $1-\alpha$. Fix the calibration set; the
coverage you then realize on fresh test points is a random variable, and it is not $1-\alpha$ on the
nose. Write the scores on the probability-integral scale, $U_i=F(S_i)$, so each $U_i$ is marginally
uniform whatever the model. With $k=\lceil(n+1)(1-\alpha)\rceil$ the conformal threshold is the order
statistic $U_{(k)}$, and the calibration-conditional coverage against an independent test draw from
the marginal is
\begin{equation}
c=U_{(k)} .
\end{equation}
For independent (or merely exchangeable, de facto independent) calibration scores this is the
classical result \citep{vovk2012,marques2024}: $c\sim\mathrm{Beta}(k,n-k+1)$, with
\begin{equation}\label{eq:betavar}
\mathbb E[c]=\frac{k}{n+1},\qquad \operatorname{Var}(c)=\frac{k(n-k+1)}{(n+1)^2(n+2)}\;\approx\;
\frac{\alpha(1-\alpha)}{n}.
\end{equation}
This is the fan: the per-dataset coverage scatters around $1-\alpha$ with a spread that shrinks like
$1/n$. The question of this note is what the dependence among the calibration scores does to that
spread. The companion notes treat the orthogonal, within-dataset axis (coverage uneven across $x$
for a fixed dataset, the information gap $I(R;X)$ and its sign \citealp{cotton_marginally,
cotton_signed}); here every statement is about the across-dataset fluctuation of the single number
$c$.

The whole problem passes through the count process $N(u)=\sum_{i=1}^n\mathbf 1\{U_i\le u\}$, because
$U_{(k)}=\int_0^1\mathbf 1\{N(u)<k\}\,\mathrm du$, so
\begin{equation}\label{eq:reduction}
\operatorname{Var}(c)=\int_0^1\!\!\int_0^1\operatorname{Cov}\big(\mathbf 1\{N(u)<k\},\,
\mathbf 1\{N(v)<k\}\big)\,\mathrm du\,\mathrm dv,
\end{equation}
a functional of the law of the count process alone. Two facts about $N$ drive everything.

\begin{lemma}[Count variance]\label{lem:count}
If the scores are negatively associated \citep{joagdev1983}, then for every $u$,
$\operatorname{Var}(N(u))\le n\,u(1-u)$, with equality iff the exceedance indicators are pairwise
uncorrelated.
\end{lemma}
\begin{proof}
The indicators $\mathbf 1\{U_i\le u\}$ are coordinatewise-monotone functions of single distinct
coordinates, hence negatively associated, hence pairwise non-positively correlated. So
$\operatorname{Var}(N(u))=\sum_i\operatorname{Var}+\sum_{i\ne j}\operatorname{Cov}\le
\sum_i u(1-u)=n\,u(1-u)$.
\end{proof}

\section{The fan coefficient}

Let $p=1-\alpha$ be the operating level and $\xi_i=\mathbf 1\{U_i\le p\}$, so $N(p)=\sum_i\xi_i$.
The Bahadur representation for the uniform quantile \citep{bahadur1966}, which holds for negatively
associated sequences under mild conditions, gives $U_{(k)}-p = p-N(p)/n + o_P(n^{-1/2})$, and hence

\begin{theorem}[Fan coefficient]\label{thm:coeff}
As $n\to\infty$,
\begin{equation}
\operatorname{Var}(c)=\frac{1}{n^2}\operatorname{Var}(N(p))\,(1+o(1))
=\frac{\alpha(1-\alpha)}{n}\,R\,(1+o(1)),\qquad
R=1+\frac{n-1}{\alpha(1-\alpha)}\,\bar c,
\end{equation}
where $\bar c=\tfrac{1}{n(n-1)}\sum_{i\ne j}\operatorname{Cov}(\xi_i,\xi_j)$ is the average pairwise
covariance of the exceedance indicators at the operating quantile.
\end{theorem}

The fan's leading coefficient is, exactly, the variance of the exceedance count at $1-\alpha$. Under
independence $\bar c=0$ and $R=1$, recovering \eqref{eq:betavar}. Negative association makes
$\bar c<0$ and $R<1$: the fan narrows. At the contest floor $\bar c=-\alpha(1-\alpha)/(n-1)$ one gets
$R=0$. The reduction factor $R$ is the single number to carry away.

\section{Positive dependence widens the fan, exactly}

When the scores are extendable, that is, drawn from an infinitely exchangeable sequence, de Finetti
\citep{definetti1937,kerns2006} makes them conditionally independent given a latent directing measure
$Q$. Conditional on $Q$ the coverage $c$ is an ordinary independent-sample order statistic, so the
law of total variance gives an exact decomposition.

\begin{theorem}[Positive dependence]\label{thm:pos}
For an extendable exchangeable calibration sequence with directing measure $Q$,
\begin{equation}
\operatorname{Var}(c)=\underbrace{\mathbb E_Q\!\big[\operatorname{Var}(c\mid Q)\big]}
_{\text{average within-dataset fan}}
+\underbrace{\operatorname{Var}_Q\!\big(\mathbb E[c\mid Q]\big)}_{\text{between-dataset term}}.
\end{equation}
The between-dataset term is non-negative and vanishes under independence. Positive cross-sample
dependence is exactly this term: a per-dataset latent that shifts the conditional quantile, on top of
the ordinary fan.
\end{theorem}

The decomposition is exact and needs no asymptotics. It says the danger of positive dependence is
not bias but variance: you draw a dataset, the latent $Q$ places the whole calibration sample high or
low, and the realized coverage rides with it. Numerically, a one-factor Gaussian model at pairwise
correlation $0.2$ has $\operatorname{Var}(c)$ about $3.3$ times the independent fan, of which the
between-dataset term is more than half.

\section{Negative dependence narrows the fan}

The negative regime is the one a genuine prior cannot reach. By \citet{kerns2006} a finite
exchangeable sequence with $\rho<0$ has a \emph{signed} directing measure, so the clean conditioning
of Theorem~\ref{thm:pos} is unavailable: there is no $Q$ to condition on, and the between-dataset
term would be a negative quasi-variance. That is precisely the regime where the fan contracts.

\begin{theorem}[Endpoints]\label{thm:ends}
Independent scores give $c\sim\mathrm{Beta}(k,n-k+1)$, the widest fan at fixed marginals. At the
maximally negatively associated floor, where the calibration scores are a uniformly random
permutation of the grid $\{1/(n+1),\dots,n/(n+1)\}$ (sampling without replacement, pairwise
correlation $-1/(n-1)$), the order statistic is deterministic, $U_{(k)}=k/(n+1)$ almost surely, so
$\operatorname{Var}(c)=0$: the realized coverage equals the nominal level on every draw.
\end{theorem}

Both endpoints agree with Theorem~\ref{thm:coeff} at $R=1$ and $R=0$. Between them the fan contracts
monotonically as the scores become more negatively dependent, and the contraction is exactly the
$R<1$ of the fan coefficient. We can prove this contraction exactly at the endpoints and to leading
order everywhere; the exact finite-sample inequality is a statement about the spread of an order
statistic under negative dependence.

\begin{conjecture}[Finite-sample contraction]\label{conj:disp}
If the calibration scores are negatively associated then $\operatorname{Var}(c)\le$ the independent
value \eqref{eq:betavar}; equivalently $U_{(k)}$ under negative association is smaller in the
dispersive order than the independent $U_{(k)}$.
\end{conjecture}

The obstruction to a one-line proof is that $\operatorname{Var}(c)$ depends on the law of $U_{(k)}$,
which negative association reshapes through $u\mapsto\mathbb P(N(u)\ge k)$; negative association makes
$N(u)$ a mean-preserving contraction of the binomial \citep{shao2000}, thinning both tails, but the
mean of $U_{(k)}$ itself shifts, so the comparison is one of spread (the dispersive order
\citealp{shaked2007}), not of a mean-matched convex order. The dispersive ordering of order
statistics under dependence is the right machinery; the natural sufficient condition is negative
supermodular dependence \citep{christofides2004}, which all the structures below satisfy. We verify
the conjecture across dependence structures (one run each, reproduced by \texttt{check\_fan.py};
$n=20$, $\alpha=0.1$, independent fan variance $3.92\times10^{-3}$).

\begin{center}
\begin{tabular}{lcccc}
\hline
calibration dependence & $\mathbb E[c]$ & $\operatorname{Var}(c)$ & ratio to fan & sign \\
\hline
independent ($\rho=0$)                 & $0.904$ & $3.92\times10^{-3}$ & $1.00$ & --- \\
exchangeable, $\rho=-1/(n-1)$          & $0.923$ & $1.16\times10^{-3}$ & $0.29$ & neg.\ assoc. \\
MA(1), lag-1 corr $-\tfrac12$          & $0.912$ & $2.71\times10^{-3}$ & $0.69$ & neg.\ assoc. \\
random negatively-paired blocks        & $0.909$ & $3.14\times10^{-3}$ & $0.80$ & neg.\ assoc. \\
contest / without replacement          & $0.905$ & $0$                  & $0.00$ & floor \\
one-factor, $\rho=+0.2$                & $0.863$ & $1.33\times10^{-2}$ & $3.41$ & positive \\
\hline
\end{tabular}
\end{center}

Every negatively associated structure, exchangeable or not, comes in under the independent fan; the
positive control widens it. The mean drifts up slightly under negative dependence and down under
positive dependence, but stays near $1-\alpha$; the action is in the variance.

\begin{remark}[The signed measure as a negative between-term]
Theorem~\ref{thm:pos} reads the fan as within-dataset variance plus a non-negative between-dataset
variance. In the signed-de-Finetti picture of the companion note, the negative corner is exactly
where that between-dataset object turns negative, a Feynman--Wigner quasi-variance that subtracts from
the within-dataset fan rather than adding to it, reaching the full cancellation $R=0$ at the contest
floor. This is heuristic, since a signed measure has no honest conditional variance, but it is the
right cartoon: the same sign that decides conditional adaptivity in \citet{cotton_signed} decides
whether the fan is inflated or cancelled here.
\end{remark}

\section{A constructive corollary}

If negative dependence narrows the fan, one can engineer it. A balanced calibration set whose scores
behave like the grid of Theorem~\ref{thm:ends} concentrates the realized coverage far inside the
independent fan, so the same reliability is bought with fewer calibration points, or the same number
of points gives less run-to-run wobble. The classical variance-reduction designs all apply in
principle: stratification never increases variance \citep{cochran1977}, antithetic pairing helps for
the monotone empirical CDF, and determinantal or repulsive designs give negative association with
faster-than-independent concentration \citep{hough2006,bardenet2016}.

\begin{remark}[The exchangeability constraint]
The marginal guarantee needs each calibration score exchangeable with the test score; it does not
need the calibration scores independent of one another. So negative association is admissible, but
only if it is induced by a map symmetric across all $n+1$ points, a pooled de-meaning or ranking,
which preserves the exchangeability of the joint. Selecting or curating the calibration set by its
scores, the obvious way to balance it, breaks exchangeability with the test point and voids the
guarantee. The viable construction is transductive and symmetric; that is also why the contest grid,
a symmetric function of the pooled sample, both keeps the marginal guarantee and zeroes the fan.
\end{remark}

\section{Related work}

The $\mathrm{Beta}(k,n-k+1)$ law of calibration-conditional coverage and its variance are
\citet{vovk2012} and \citet{marques2024}; the entire prior literature treats the calibration scores
as independent or exchangeable and reads off a single number, never varying the dependence to move
it. Negative association and its consequences are \citet{joagdev1983}; that negatively associated
sums are dominated in convex order by independent ones is \citet{shao2000}, and sampling without
replacement concentrates at least as tightly as with replacement \citep{serfling1974}. The dispersive
order and its variance consequence are standard \citep{shaked2007}; the connection between supermodular
ordering and negative dependence is \citet{christofides2004}; the Bahadur representation is
\citet{bahadur1966}. Determinantal and repulsive designs are \citet{hough2006,bardenet2016} and
stratification is \citet{cochran1977}. The finite signed de Finetti representation is
\citet{kerns2006}, and the within-dataset, across-$x$ reading of its sign is the companion note
\citep{cotton_signed}; the information gap is \citep{cotton_marginally}. The contribution here is to
make the dependence the control variable for the \emph{variance} of realized coverage: the exact
fan coefficient $R$, the exact de Finetti decomposition for the positive side, the exact zero at the
contest floor, and the reduction of the general negative case to a dispersive ordering.

\begin{thebibliography}{99}
\bibitem[Bahadur(1966)]{bahadur1966} Bahadur, R.~R. (1966). A note on quantiles in large samples.
\emph{The Annals of Mathematical Statistics} 37(3):577--580.
\bibitem[Bardenet and Hardy(2020)]{bardenet2016} Bardenet, R. and Hardy, A. (2020). Monte Carlo with
determinantal point processes. \emph{The Annals of Applied Probability} 30(1):368--417.
arXiv:1605.00361.
\bibitem[Christofides and Vaggelatou(2004)]{christofides2004} Christofides, T.~C. and Vaggelatou, E.
(2004). A connection between supermodular ordering and positive/negative association. \emph{Journal
of Multivariate Analysis} 88(1):138--151.
\bibitem[Cochran(1977)]{cochran1977} Cochran, W.~G. (1977). \emph{Sampling Techniques}, 3rd ed.
Wiley.
\bibitem[Cotton()]{cotton_marginally} Cotton, P. Marginally Useful: Formalizing the Information Gap
in Conformal Prediction. Companion paper.
\bibitem[Cotton()]{cotton_signed} Cotton, P. A Feynman--Wigner-Style Diagnostic for the Efficacy of
Conformal Prediction via Signed de Finetti Representations. Companion note.
\bibitem[de Finetti(1937)]{definetti1937} de Finetti, B. (1937). La pr\'evision: ses lois logiques,
ses sources subjectives. \emph{Annales de l'Institut Henri Poincar\'e} 7(1):1--68.
\bibitem[Hough et~al.(2006)]{hough2006} Hough, J.~B., Krishnapur, M., Peres, Y. and Vir\'ag, B.
(2006). Determinantal processes and independence. \emph{Probability Surveys} 3:206--229.
\bibitem[Joag-Dev and Proschan(1983)]{joagdev1983} Joag-Dev, K. and Proschan, F. (1983). Negative
association of random variables, with applications. \emph{The Annals of Statistics} 11(1):286--295.
\bibitem[Kerns and Sz\'ekely(2006)]{kerns2006} Kerns, G.~J. and Sz\'ekely, G.~J. (2006). De Finetti's
theorem for abstract finite exchangeable sequences. \emph{Journal of Theoretical Probability}
19(3):589--608.
\bibitem[Marques F.(2024)]{marques2024} Marques F., P.~C. (2024). Universal distribution of the
empirical coverage in split conformal prediction. \emph{Statistics \& Probability Letters} 219:110350
(2025). arXiv:2303.02770.
\bibitem[Serfling(1974)]{serfling1974} Serfling, R.~J. (1974). Probability inequalities for the sum
in sampling without replacement. \emph{The Annals of Statistics} 2(1):39--48.
\bibitem[Shaked and Shanthikumar(2007)]{shaked2007} Shaked, M. and Shanthikumar, J.~G. (2007).
\emph{Stochastic Orders}. Springer.
\bibitem[Shao(2000)]{shao2000} Shao, Q.-M. (2000). A comparison theorem on moment inequalities
between negatively associated and independent random variables. \emph{Journal of Theoretical
Probability} 13(2):343--356.
\bibitem[Vovk(2012)]{vovk2012} Vovk, V. (2012). Conditional validity of inductive conformal
predictors. \emph{Asian Conference on Machine Learning}, PMLR 25:475--490.
\end{thebibliography}

\end{document}
